\documentclass[DIN, pagenumber=false, fontsize=11pt, parskip=half]{scrartcl}
\usepackage[UTF8]{ctex}  % 中文支持
% \usepackage{ngerman}
\usepackage[utf8]{inputenc}
\usepackage[T1]{fontenc}
\usepackage{textcomp}
\usepackage{xcolor} % Required for custom red

\usepackage{geometry} % Custom margins

\usepackage{tocloft} % Modifying the table of contents

% for matlab code
% bw = blackwhite - optimized for print, otherwise source is colored
\usepackage[framed,numbered,bw]{mcode}
\usepackage[most]{tcolorbox}
% for other code
%\usepackage{listings}

\setlength{\parindent}{0em}

% set section in CM
\setkomafont{section}{\normalfont\bfseries\Large}

\newcommand{\mytitle}[1]{{\noindent\Large\textbf{#1}}}
\newcommand{\mysection}[1]{\textbf{\section*{#1}}}
\newcommand{\mysubsection}[2]{\romannumeral #1) #2}

\newtcolorbox{parchmentquote}{
    enhanced,
    colback=orange!5!white,   % 浅米黄色背景
    colframe=orange!50!black,
    boxrule=0.5pt,
    arc=2pt,                  % 细微圆角
    boxsep=5pt,
    before skip=10pt,
    after skip=10pt,
    fontupper=\itshape,
    borderline west={2pt}{0pt}{orange!80!black} % 加厚的左装饰线
}
%===================================
\begin{document}

\noindent\textbf{课下讲义} \hfill \textbf{online course}\\
primary school \hfill Liu\\

\mytitle{行程问题 \hfill \today}
\newline
\begin{parchmentquote}
	行程为题是数学与物理学中常见的题型，涉及到速度、时间和距离之间的关系。难一点的题目需要借助图像法，也就是柳卡图等来解答。
\end{parchmentquote}


%===================================
\section{简单行程问题}
$ s = v \times t $
\newline
其中，\(s\) 表示距离，\(v\) 表示速度，\(t\) 表示时间。
\newline
\subsection{平均速度} 
已知上山的平均速度为30米/分，上山用时85分钟，下山原路返回用时51分钟，求下山的平均速度。

\subsection{水中行船} 
顺水船速 = 船速 + 水速

逆水船速 = 船速 - 水速
\newline
一艘船往返于相距648千米的两个码头之间，顺水而下需要36个小时，逆流而下需要54个小时。
如今，水流速度增加，顺水而行只需要27小时，那么现在逆流而上需要多久。
\section{相遇问题}
$s = (v_1 + v_2) \times t$,
\subsection{同向出发相遇}
姐弟俩同时从家里出发去图书馆，路程全长1540米，弟弟每分钟走70米，姐姐骑自行车忆每分钟150米的速度到达图书馆后立即返回，途中与弟弟相遇。
这时弟弟走了几分钟？
\subsection{多人多物相遇}
甲乙两人相距30km，甲速度为$3km/h$，乙速度为$7km/h$，在某时刻，两人同向而行。甲带了一个机器人，机器人以$10km/h$的速度在甲和乙之间来回行走。问：当甲和乙相遇时，机器人共走了多少公里？
\subsection{圆形跑道}
在长为400米的环形跑道上，甲和乙从同一个地点开始相向跑步。甲的速度是每分钟120米，乙的速度是每分钟80米。两人相遇后继续跑，从出发开始，经过多久两人第二次相遇？
\end{document}